\scanpage{14}
\mspage{2}
Mais on ne peut pas dire \inserted{en général} que $\dot T$ soit la frontière
de $\bar T$ dans $X$, i.e. (\unsure{7|8}) devient faux si $\lambda=1$
(prendre $X=\bar T\deleted{\unsure{_\lambda}}$, avec $\dot T\neq0$).

\deleted{On dira} \unsure{que} dire qu'une partie fermée $\unsure{Y|Z}$ d'un
espace topologique $X$ est une ``\uline{partie bord}''
s'il existe un \uline{voisinage} \deleted{\uline{collier}}
\inserted{\uline{cylindrique}} de $\unsure{Y|Z}$ dans $X$,
i.e. un voisinage \deleted{tubulaire} conique $\bar T$ tel que
$\dot T\longrightarrow Y$ soit un homéom. --- ou ce qui revient
au même, \deleted{i.e. tel que l'on} s'il existe un
voisinage $\deleted{\unsure{de}}\bar T$ de $Y$ dans $X$ et un homéom.
\[
(9)\qquad\bar T\simeq Y\times I
\]
rendant commutatif le diagramme
\diagram{\[
(10)\qquad
\begin{tikzcd}
Y\arrow[r,hook]\arrow[dr,"{y\mapsto(y,0)}"']&\bar T\arrow[d,"\simeq",draw=none]\\
&Y\times I
\end{tikzcd}
\]}
\margin{définir la \unsure{structure} de voisinage cylindrique
comme un homéom.
$E\simeq Y\times[0,1]$
compatible avec
\diagram{\[
\begin{tikzcd}
Y\arrow[r]&\unsure{[0,1]|Y\times[0,1]}\\
y\arrow[r,mapsto]&(y,0)
\end{tikzcd}
\]}}

Par exemple, si $\bar T$ est un voisinage \deleted{tubulaire} conique
\inserted{(absolu)} de \deleted{\unsure{tout}} $Y$, alors $\dot T$ est un sous-espace bord.
\unsure{donc} $\unsure{\bar T_\lambda|\dot T_\lambda}$ est un sous-espace bord de $\bar T_\lambda$ ;
$\deleted{\unsure{\mathrm{Ep}}}\dot T,(1-\unsure{\varepsilon|\lambda})\dot T$
et $\dot T\cup(1-\varepsilon)\dot T$ sont des sous-espaces
bords de $\mathrm{Ep}_\varepsilon^{(T)}$. \deleted{Dans}
Si $X$ est un espace
\uline{compact}, et $(Y_i)_{i\in I}$ une \unsure{partie|famille} finie de
parties fermées disjointes de $X$, alors
$Y=\bigcup_{i\in I}Y_i$ est un sous-espace bord ssi les $Y_i$ le sont.
\margin{13}
